%----------------------------------------------------------------------------------------
% NOTATIE EN FIGUREN VAN DE CURSUS STATISTIEK (2de graad)
% Staat in de preambule en niet in een hoofdstuk: ook de lijsten met oplossingen gebruiken dit.
% Vertrokken van notatie.tex van Integralen; wat hier niet nodig is, is geschrapt.
%----------------------------------------------------------------------------------------
% Bewaarde grafiek, geplaatst door maak-pdf.py: bestand en pagina in figuren-cache/
\newcommand{\grafiekbestand}[2]{\includegraphics[page=#2,trim=3bp 3bp 3bp 3bp]{figuren-cache/#1}}

% Statistische notatie
\providecommand{\xbar}{\overline{x}}                  % gemiddelde
\newcommand{\Me}{\mathrm{Me}}                         % mediaan
\newcommand{\Mo}{\mathrm{Mo}}                         % modus
\newcommand{\IKA}{\mathrm{IKA}}                       % interkwartielafstand
\newcommand{\VB}{\mathrm{VB}}                         % variatiebreedte
\providecommand{\abs}[1]{\lvert #1\rvert}

% Zoals stappen, maar met een bredere kolom voor de berekening: voor lange formules
\newenvironment{stappenbreed}{%
  \par\renewcommand{\arraystretch}{1.25}%
  \noindent\begin{tabular}{@{}>{$\displaystyle}p{9.8cm}<{$}@{\hspace{4mm}}>{\RaggedRight\small\itshape\arraybackslash}p{4.9cm}@{}}%
}{\end{tabular}\par}

% Zoals stappen, met eigen kolombreedtes: \begin{stappenvrij}{6.2cm}{3.8cm}. Voor naast een figuur.
\newenvironment{stappenvrij}[2]{%
  \par\renewcommand{\arraystretch}{1.25}%
  \noindent\begin{tabular}{@{}>{$\displaystyle}p{#1}<{$}@{\hspace{3mm}}>{\RaggedRight\small\itshape\arraybackslash}p{#2}@{}}%
}{\end{tabular}\par}

% Tekst of berekening links, figuur rechts: \naast{links}{rechts}
\newsavebox{\naastbox}
\newcommand{\naast}[3][]{%
  \par\noindent\sbox{\naastbox}{#3}%
  \begin{minipage}[t]{\dimexpr\linewidth-\wd\naastbox-3mm\relax}\vspace{0pt}#2\end{minipage}\hfill
  \begin{minipage}[t]{\wd\naastbox}\vspace{0pt}\usebox{\naastbox}\end{minipage}\par\addvspace{1mm}}
% Een rij figuren, gecentreerd
\newcommand{\figuren}[1]{\par\addvspace{1mm}{\centering #1\par}\addvspace{1.5mm}}
% Opgave met een figuur rechts: \begin{metgrafiek}{tikz} opgave en deelvragen \end{metgrafiek}
\newsavebox{\grafiekbox}
\newenvironment{metgrafiek}[2][]{%
  \sbox{\grafiekbox}{#2}%
  \begin{minipage}[t]{\dimexpr\linewidth-\wd\grafiekbox-3mm\relax}%
}{%
  \end{minipage}\hfill
  \begin{minipage}[t]{\wd\grafiekbox}\vspace{-2mm}\usebox{\grafiekbox}\end{minipage}\par}

% Een tabel met gegevens, klein en gecentreerd: \gegevens{c|cccc}{rijen}
\newcommand{\gegevens}[2]{\figuren{\small\renewcommand{\arraystretch}{1.2}\setlength{\tabcolsep}{4.5pt}\begin{tabular}{#1}#2\end{tabular}}}

%----------------------------------------------------------------------------------------
% DIAGRAMMEN, binnen een tikzpicture. Zet de schaal met x=... en y=...
%----------------------------------------------------------------------------------------
% Horizontale getallenas op hoogte 0: \getalas[eenheid]{min}{stap}{max}, met streepjes en getallen
\newcommand{\getalas}[4][]{%
  \draw[->,grijs] (#2,0) -- ({#4+.5*(#3)},0) node[right,text=black,font=\scriptsize] {#1};
  \pgfmathsetmacro{\gatwee}{#2+#3}%
  \foreach \t in {#2,\gatwee,...,#4}{\draw[grijs] (\t,-.07) -- (\t,.07);
    \node[below,font=\scriptsize,inner sep=2pt,text=black,text height=1.4ex,text depth=.2ex] at (\t,-.07) {\pgfmathprintnumber[use comma,1000 sep={}]{\t}};}}
% Verticale as met lichte rasterlijnen: \frequentieas[titel]{stap}{max}{xmin}{xmax}
\newcommand{\frequentieas}[5][frequentie]{%
  \pgfmathsetmacro{\fatwee}{#2+#2}%
  \foreach \t in {#2,\fatwee,...,#3}{\draw[black!12] (#4,\t) -- (#5,\t);
    \node[left,font=\scriptsize,inner sep=2pt,text=black] at (#4,\t) {\pgfmathprintnumber[use comma,1000 sep={}]{\t}};}
  \draw[->,grijs] (#4,0) -- (#4,{#3+.6*(#2)}) node[above,text=black,font=\scriptsize] {#1};
  \draw[grijs] (#4,0) -- (#5,0);}
% Staaf in een staafdiagram: \staaf{midden}{hoogte}{breedte}
\newcommand{\staaf}[3]{\fill[accent!40] ({#1-#3/2},0) rectangle ({#1+#3/2},#2); \draw[accent] ({#1-#3/2},0) rectangle ({#1+#3/2},#2);}
% Balk in een histogram: \balk{van}{tot}{hoogte}. De balken raken elkaar.
\newcommand{\balk}[3]{\fill[accent!30] (#1,0) rectangle (#2,#3); \draw[accent] (#1,0) rectangle (#2,#3);}
% Boxplot op hoogte y: \boxplot[hoogte]{y}{min}{Q1}{Me}{Q3}{max}
\newcommand{\boxplot}[7][.6]{%
  \draw[accent,thick] (#3,#2) -- (#4,#2) (#6,#2) -- (#7,#2);
  \draw[accent,thick] (#3,{#2-#1/3}) -- (#3,{#2+#1/3}) (#7,{#2-#1/3}) -- (#7,{#2+#1/3});
  \fill[accent!18] (#4,{#2-#1/2}) rectangle (#6,{#2+#1/2});
  \draw[accent,thick] (#4,{#2-#1/2}) rectangle (#6,{#2+#1/2});
  \draw[accent,line width=1.8pt] (#5,{#2-#1/2}) -- (#5,{#2+#1/2});}
% Uitschieter bij een boxplot: \uitschieter{x}{y}
\newcommand{\uitschieter}[2]{\draw[accent,thick] (#1,#2) circle (2pt);}
% Stip in een stippendiagram: \stip{x}{hoeveelste}
\newcommand{\stip}[2]{\fill[accent] (#1,{.18+.22*(#2-1)}) circle (2.2pt);}
% Sector in een cirkeldiagram: \sector{van graden}{tot graden}{kleur}{label}
\newcommand{\sector}[4]{%
  \fill[#3] (0,0) -- (#1:1) arc[start angle=#1,end angle=#2,radius=1] -- cycle;
  \draw[white,thick] (0,0) -- (#1:1) arc[start angle=#1,end angle=#2,radius=1] -- cycle;
  \node[font=\scriptsize] at ({(#1+#2)/2}:1.38) {#4};}
% Turven: \turf{n} tekent n streepjes, per vijf gebundeld met een schuine streep
\newcommand{\turf}[1]{\begin{tikzpicture}[baseline=-.4mm,x=.85mm,y=2.6mm,line width=.55pt]
  \pgfmathtruncatemacro{\turfv}{int(#1/5)}\pgfmathtruncatemacro{\turfr}{#1-5*\turfv}%
  \ifnum\turfv>0 \foreach \g in {1,...,\turfv}{\foreach \s in {0,...,3}{\draw ({(\g-1)*6+\s},0) -- ++(0,1);}
    \draw ({(\g-1)*6-.6},.2) -- ({(\g-1)*6+3.6},.8);}\fi
  \ifnum\turfr>0 \foreach \s in {1,...,\turfr}{\draw ({\turfv*6+\s-1},0) -- ++(0,1);}\fi
\end{tikzpicture}}
